% `tex.web` §143/§144 and §193: a ligature keeps the characters it replaced.
%
% When the font's ligature program substitutes a character for a run, TeX does
% not simply append the result — it records the run as well, and §193 displays
% both:
%
%     .\r ^^L (ligature fi)
%
% The character actually set, then what it came from in §174's short form.
%
% We applied the program and appended a plain character. The DVI was identical,
% which is why tools/typeset_check.sh (30/30) and tools/latex_check.sh (35/35)
% never noticed: the glyph reaching the page is the same one. Found in trip.tex's
% sixth shipped page, where Knuth's transcript has `\rip C (ligature B-)`.
%
% It is not only a display difference. §897's `reconstitute` takes a word apart
% to hyphenate it and re-runs the ligature program on the letters; `\unhbox`
% hands the nodes back; §1124's `\accent` inspects what precedes it. A document
% that *hyphenates* a word containing a ligature is the case at risk, not one
% that merely sets one.
%
% Also §48/§49: a character outside ` `..`~` cannot be printed, so TeX writes
% `^^X` — the code plus 64 below 64, minus 64 above it. Ligature characters are
% usually in that range (cmr10's `fi` is 12), which is why it surfaced here.
\catcode`\{=1 \catcode`\}=2
\showboxdepth=2 \showboxbreadth=9 \tracingonline=0
\font\r=cmr10 \r

% fi is one step; ffi is two, and the run accumulates across both
\setbox0\hbox{fi ffi --- A}
\showbox0

% a ligature that is itself printable needs no ^^ notation
\setbox1\hbox{---}
\showbox1

% and a word with no ligature in it is unchanged
\setbox2\hbox{box}
\showbox2
\end
